% RecSys Odyssey repository-backed lecture note.
% Add figures with: \coursefigure{figures/name.svg}{Caption}
\section{Recall and coverage answer different questions}

\begin{abstract}
Recall checks whether good items survived; coverage checks whether the system has breadth.
\end{abstract}

\subsection{Concept}
Recall@k asks what fraction of known relevant items reached the top-k. It is especially important for candidate generation: a ranker cannot rescue an item that retrieval dropped. Coverage measures how much of the catalogue, categories or user population receives recommendations. Neither guarantees a satisfying slate, so teams read them beside NDCG, diversity, latency and online outcomes.

\begin{equation*}
Recall@k = |relevant ∩ top-k| / |relevant|      item coverage = |recommended items| / |catalogue|
\end{equation*}

\subsection{Key terms}
\begin{itemize}
\item \textbf{Recall@k.} The fraction of relevant items present in the first k results.
\item \textbf{Coverage.} The breadth of items, categories or users reached by a policy.
\item \textbf{Diversity.} How different the items within one slate are from each other.
\end{itemize}
